\documentclass[tikz,border=6pt]{standalone}
\usepackage{peabra}
\begin{document}
\begin{tikzpicture}[x=1cm,y=1cm,
  st/.style={draw=peVilagos, fill=peHalvany, text=peSotet, rounded corners=1.5pt,
    font=\scriptsize, align=center, inner sep=2pt, minimum height=9mm, text width=17mm, line width=0.6pt},
  resz/.style={rounded corners=2pt, font=\scriptsize, align=left, inner sep=3pt, text=white, anchor=west},
  fej/.style={font=\scriptsize, text=peSotet, align=left, anchor=west}]
  \foreach \n/\t in {1/{Rendszer és\\környezet},2/{Veszély-\\azonosítás},3/{Kockázat-\\értékelés},
      4/{Biztonsági\\követelmények},5/{Tervezés,\\architektúra},6/{Kvantitatív\\igazolás},
      7/{Verifikáció,\\érvelés},8/{Üzemeltetés,\\tanulás}}
    \node[st] (s\n) at ({(\n-1)*1.98},0) {\t};
  \foreach \a/\b in {1/2,2/3,3/4,4/5,5/6,6/7,7/8}
    \draw[-{Stealth[length=1.3mm]}, draw=peSzurke, line width=0.5pt] (s\a) -- (s\b);
  % részek sávjai: kezdő állomás, záró állomás, magasság, szín, felirat
  \def\sav#1#2#3#4#5{%
    \fill[#4, rounded corners=2pt] ($(s#1.west)+(0,#3)$) rectangle ($(s#2.east)+(0,#3-0.62)$);
    \node[resz] at ($(s#1.west)+(0.05,#3-0.31)$) {#5};}
  \sav{1}{8}{-0.85}{peSotet}{\textbf{I. Alapok} \quad 1. balesetek és balesetmodellek \textperiodcentered{} 2. fogalmak \textperiodcentered{} 3. kockázatértékelés \textperiodcentered{} 4. életciklus és szabványok}
  \sav{2}{3}{-1.6}{peKek}{\textbf{II. Veszélyazonosítás}\quad 5--7.}
  \sav{3}{7}{-2.35}{peKozep}{\textbf{III. Kvantitatív megbízhatóság és biztonság}\quad 8--13.}
  \sav{4}{8}{-3.1}{peZold}{\textbf{IV. Tervezés biztonságra}\quad 14. architektúrák \textperiodcentered{} 15. szoftver \textperiodcentered{} 16. ember}
  \sav{6}{8}{-3.85}{peLila}{\textbf{V. Bizonyítás, új kihívások}\quad 17--18.}
  % esetek
  \node[fej] at ($(s1.west)+(0,-4.85)$) {\textbf{Futó esetek:}\quad
    {\color{esetA}\rule{2.4mm}{2.4mm}}~A -- ReaktorVéd (folyamatipar)\qquad
    {\color{esetB}\rule{2.4mm}{2.4mm}}~B -- VészFék (közúti jármű)\qquad
    {\color{esetC}\rule{2.4mm}{2.4mm}}~C -- KoboCella (gépipar, robotika)};
\end{tikzpicture}
\end{document}
